\section{Exceptional Local Profiles and Vd Placements}
\label{sec:short-vd-placements}

% STATEMENTS-ONLY-REPLACE-BEGIN: vd-profile-introduction
This subsection proves the local profile estimates used by the
one-special-triangle branches.  A T3-like adjacent exceptional triangle and a
Vd1 adjacent exceptional triangle satisfy the same two scalar
inequalities and therefore share one center-hiding theorem.  The remaining
Vd1/Vd2 placements use the interval residuals from
Section~\ref{sec:universal-calculus}, one global radial envelope, and two
immediate consequences of the corner normal form.
% STATEMENTS-ONLY-REPLACE-END: vd-profile-introduction

Use the strict-supercritical envelope $M_c^{\rm sup}$ from
\eqref{eq:reader-free-supercritical-envelope}; its explicit formula is
strictly decreasing for $0\le c<1/2$.  At $c=1/2$ we use only the continuous
extension $M_{1/2}^{\rm sup}=1/2$.  The strict feasible set is empty there,
so this endpoint value is never described as an attained supremum.

\subsection{A global quarter radial envelope}

For $(p,h)\in[0,1]^2$, let $c_{\max}(p,h)$ be the maximal radial coordinate
such that $(p,h,c)$ belongs to the local admissible set.

\begin{lemma}[Quarter radial envelope]
\label{lem:quarter-radial-envelope}
If
\[
 0\le h\le p,
 \qquad
 p+h<1,
\]
then
\begin{equation}
 c_{\max}(p,h)\le1-\frac h4.
 \label{eq:quarter-radial-envelope}
\end{equation}
The inequality is strict on the selected $\mathcal A_T$ component.
\end{lemma}

\begin{proof}
Put $s=p+h$ and $q=s^4-s^2+ph$.  Since $s<1$, only the two
nonsupercritical cells of Proposition~\ref{prop:exact-admissible-set} occur.

On $\mathcal A_L$, the smaller boundary coordinate is $h$, and the selected
radial frontier is the unique root in $[\sqrt3/2,1]$ of
\[
 P_h(c)=c^4-c^2+hc-h^2=0.
\]
At $c_0=1-h/4$,
\[
 P_h(c_0)=\frac h{256}(h^3-16h^2-240h+128).
\]
The cubic decreases on $[0,1/2]$ and has value $33/8$ at $1/2$, so this is
positive for $h>0$.  Moreover
\[
 P_h\left(\frac{\sqrt3}{2}\right)
 =-\left(h-\frac{\sqrt3}{4}\right)^2\le0,
 \qquad
 P_h'(c)>0\quad(c\ge\sqrt3/2).
\]
Hence the selected root is below $c_0$; for $h=0$ equality gives
$c_{\max}=1$.

On $\mathcal A_T$, put
\[
 r=\sqrt{4s^2-3},
 \qquad
 \gamma=\frac{p-h}{s}.
\]
The exact identity
\[
 \frac q{s^2}=\frac{r^2-\gamma^2}{4}
\]
gives $0\le\gamma<r$.  Write $\gamma=wr$, $0\le w<1$.  Then
\[
 p=\frac{s(1+wr)}2,
 \qquad
 h=\frac{s(1-wr)}2,
 \qquad
 c_{\max}=\frac{s(1+wr)}{1+r}.
\]
For
\[
 \Psi(w)=c_{\max}+\frac h4
\]
one has
\[
 \Psi'(w)=\frac{sr(7-r)}{8(1+r)}>0,
\]
so
\[
 \Psi(w)<\Psi(1)=\frac{s(9-r)}8.
\]
Since $4s^2=r^2+3$,
\[
 256-(r^2+3)(9-r)^2
 =(1-r)(r^3-17r^2+67r+13)>0
\]
for $0\le r<1$; the cubic is increasing and positive on this interval.
Thus $s(9-r)<8$ and $c_{\max}+h/4<1$.
\end{proof}

\subsection{The rational T3-like adjacent-exception profile}

\begin{lemma}[T3-like adjacent-exception profile]
\label{lem:short-t3-rescuer-profile}
Let a translated T3-like triangle at $V_0$ contain $M_1$, have boundary trace
$[0,a]$ on $e_{5,0}$, and have trace $[c,u]$ on $r_1$ measured
from $V_1$ toward $O$.  Then
\begin{equation}
 a\le 1-M_c^{\rm sup},
 \qquad
 \frac{a}{a+1-u}\le 1-M_c^{\rm sup}.
 \label{eq:short-t3-profile}
\end{equation}
\end{lemma}

\begin{proof}
The translated T3-like normal form supplies
\[
 1\le D\le\frac2{\sqrt3},
 \qquad
 E=\sqrt{4-3D^2},
 \qquad
 \varrho=\frac{D+E}{2},
\]
and
\[
 \frac{E}{2D}\le a\le\frac{4-3D+E}{4D}.
\]
Its radial interval is
\[
 c=\frac{D(1+a)-1}{\varrho},
 \qquad
 u=1-\frac\varrho D+a.
\]
Put
\[
 x=\frac{aD}{\varrho},
 \qquad
 \theta=\frac{D-1}{\varrho}.
\]
The relation $\varrho^2-D\varrho+D^2=1$ gives
\[
 \varrho=\frac{1-2\theta}{1-\theta+\theta^2},
 \qquad
 D=\frac{1-\theta^2}{1-\theta+\theta^2},
 \qquad
 E=\frac{1-4\theta+\theta^2}{1-\theta+\theta^2},
\]
with $0\le\theta\le2-\sqrt3$.  Moreover
\begin{equation}
 c=x+\theta,
 \qquad
 \frac{1-4\theta+\theta^2}{2(1-2\theta)}
 \le x\le\frac12-\theta.
 \label{eq:short-t3-x-domain}
\end{equation}

Let
\[
 \psi(z)=\frac{z(2-z)}{1+z}.
\]
It is increasing on $[0,1/2]$ and
$\psi(1-M_c^{\rm sup})=c$.  Hence it is enough to prove
$\psi(x)\le x+\theta$, or equivalently
\[
 Q_\theta(x)=2x^2+(\theta-1)x+\theta\ge0.
\]
The vertex is $x_*=(1-\theta)/4$.  If $0\le\theta\le1/5$, the lower endpoint
$x_L$ in \eqref{eq:short-t3-x-domain} satisfies
\[
 x_L-x_*=\frac{1-5\theta}{4(1-2\theta)}\ge0,
\]
and
\[
 Q_\theta(x_L)
 =\frac{\theta(1-5\theta+11\theta^2-\theta^3)}
 {2(1-2\theta)^2}\ge0.
\]
If $1/5\le\theta\le2-\sqrt3$, the vertex lies in the feasible interval and
\[
 Q_\theta(x)\ge Q_\theta(x_*)
 =\frac{10\theta-1-\theta^2}{8}>0.
\]
Thus $x\le 1-M_c^{\rm sup}$.  Since $D\ge\varrho$, one has $a\le x$, and
\[
 \frac{a}{a+1-u}=\frac{aD}{\varrho}=x.
\]
This proves both inequalities.
\end{proof}

\subsection{The Vd1 adjacent-exception profile}

\begin{lemma}[Vd1 adjacent-exception profile]
\label{lem:short-vd1-rescuer-profile}
Let a Vd1 triangle at $V_0$ contain $M_1$, have boundary trace $[0,a]$ on
$e_{5,0}$, and have trace $[c,u]$ on $r_1$.  Assume it has no
positive support on $r_5$.  Then
\begin{equation}
 a\le 1-M_c^{\rm sup},
 \qquad
 \frac{a}{a+1-u}\le 1-M_c^{\rm sup}.
 \label{eq:short-vd1-profile}
\end{equation}
\end{lemma}

\begin{proof}
The Vd corner normal form gives $t>0$, $d=\sqrt{t^2+t+1}$, and
\[
 c=\frac{t(1-b)}{t+1},
 \qquad
 u=\frac{d-a-tb-1}{t}.
\]
The midpoint condition implies
\[
 tb=t-c(t+1),
 \qquad
 a+tb\le d-1-\frac t2.
\]
One-sided support forces $t\ge1$.  Indeed, if $t<1$, then
\[
 d-a-tb-t\ge1-\frac t2>\frac1{t+1}>\frac{1-a}{t+1},
\]
so the raw $r_5$ interval has positive length, contrary to the Vd1 type.

Put
\[
 F(t,c)=\sqrt{t^2+t+1}-1-\frac{3t}{2}+c(t+1).
\]
Then $a\le F(t,c)$.  Since $c\le1/2$,
\[
 \frac{\partial F}{\partial t}
 =\frac{2t+1}{2\sqrt{t^2+t+1}}-\frac32+c
 \le\frac{2t+1}{2\sqrt{t^2+t+1}}-1<0,
\]
where the last inequality follows from
$(2t+1)^2<4(t^2+t+1)$.  Thus $F$ decreases for $t\ge1$, and
\[
 a\le F(t,c)\le L(c):=\sqrt3-\frac52+2c.
\]
Every feasible V triangle has $L(c)\ge0$.  The inequality
$2L(c)\le 1-M_c^{\rm sup}$ is equivalent to
\[
 \sqrt{c^2-8c+4}\le12-4\sqrt3-9c.
\]
The right side is positive on $[0,1/2]$.  Squaring is therefore
legitimate, and the squared difference is $4Q(c)$, where
\[
 Q(c)=20c^2+(18\sqrt3-52)c+47-24\sqrt3.
\]
The polynomial is decreasing on $[0,1/2]$ and
$Q(1/2)=26-15\sqrt3>0$.  Hence $2L(c)\le 1-M_c^{\rm sup}$,
and therefore $a\le 1-M_c^{\rm sup}$.

Put $\varepsilon=1-u>0$.  The endpoint formula gives
\[
 \varepsilon=\varepsilon_0+\frac at,
 \qquad
 \varepsilon_0=\frac{2t+1-c(t+1)-d}{t}.
\]
Since $c\le1/2$ and $t\ge1$,
\[
 2t+1-c(t+1)-d\ge\frac{3t+1}{2}-d>0;
\]
the last inequality follows after squaring from
$(3t+1)^2/4-d^2=(5t^2+2t-3)/4>0$.  Thus
$\varepsilon_0>0$.  Moreover,
\[
 \varepsilon_0+\frac{F(t,c)}t=\frac12.
\]
Since $z/(z+\varepsilon_0+z/t)$ increases in $z$,
\[
 \frac{a}{a+\varepsilon}
 \le\frac{F(t,c)}{F(t,c)+1/2}
 \le\frac{L(c)}{L(c)+1/2}
 \le2L(c)
 \le 1-M_c^{\rm sup}.
\]
\end{proof}

\subsection{A common adjacent-exception theorem}

\begin{lemma}[Common adjacent-exception obstruction]
\label{lem:signed-adjacent-rescuer}
Let a local adjacent exceptional triangle $T_0$ have trace $[0,a]$ on $e_{5,0}$ and
trace $[c,u]$ on $r_1$, measured from $V_1$ toward $O$.  Put
\[
 \varepsilon=1-u>0,
 \qquad
 \Theta=\frac{a}{a+\varepsilon}.
\]
Assume
\begin{equation}
 a\le 1-M_c^{\rm sup},
 \qquad
 \Theta\le 1-M_c^{\rm sup},
 \label{eq:signed-rescuer-local}
\end{equation}
and assume radial isolation forces the center to cover the $O$-side gap of
length $\varepsilon$.  If $T_1$ is the unique strict supercritical V triangle and
$T_2,T_3,T_4,T_5$ are nonsupercritical, then the perimeter and $r_1$ are not
covered.
\end{lemma}

\begin{proof}
The center exit on $r_1$ is $\delta/R$, so radial isolation gives
\[
 \delta\ge R\varepsilon.
\]
Let $h$ be the reach forced on $T_5$ along $e_{5,0}$.  If the companion center
trace does not hide the gap after coordinate $a$, then
\[
 h\ge1-a\ge M_c^{\rm sup}.
\]
In the only hiding order,
\[
 \frac{k}{W}\le a<R+\alpha.
\]
The left endpoint inequality and radial isolation give
\[
 k\le Wa,
 \qquad
 R\varepsilon\le Wa,
\]
so $R(a+\varepsilon)\le a$ and $R\le\Theta$.  Also
\[
 \alpha\le Wa-R\varepsilon-\eta,
\]
whence
\[
 R+\alpha\le a+R(u-a).
\]
The hiding order forces $u>a$, and therefore
\[
 R+\alpha\le a+\Theta(u-a)=\Theta\le 1-M_c^{\rm sup}.
\]
Thus in every ordering
\[
 h\ge M_c^{\rm sup}.
\]
Coverage of $r_1$ gives the supercritical V triangle radial lower bound at least $c$, so
its actual forward reach satisfies
\[
 B_1<M_c^{\rm sup}\le h.
\]
Lemma~\ref{lem:reader-boundary-path}, applied to $T_2,T_3,T_4,T_5$, gives
\[
 \sum_{i=2}^5(A_i+B_i)\ge4+h-B_1>4,
\]
contrary to their four nonsupercritical caps.
\end{proof}

\begin{corollary}[Adjacent T3-like/Vd1 rescue]
\label{cor:signed-adjacent-rescuers}
\label{thm:s2-geom-sc}
The common adjacent-exception lemma applies both to the exactly-one-T3-like
branch and to a normalized Vd1 triangle $T_i$ containing $M_{i-1}$ or
$M_{i+1}$.
\end{corollary}

\begin{proof}
Lemmas~\ref{lem:short-t3-rescuer-profile} and
\ref{lem:short-vd1-rescuer-profile} give the two local inequalities.  In both
branches, midpoint and support isolation force the center to cover the
$O$-side radial gap before the adjacent exceptional triangle begins.
\end{proof}

\subsection{Adjacent Vd1/Vd2 placement}

\begin{lemma}[Adjacent Vd1/Vd2 placement]
\label{lem:signed-vd-adjacent-placement}
\label{thm:s2-geom-vd-adjacent}
Assume a reduced CE2 candidate with
$T_C\cap\{M_0,\ldots,M_5\}=\{M_0\}$ has $N_+=1$, V triangle $T_0$ uniquely
supercritical, V triangle $T_1$ the unique Vd1/Vd2 V triangle, and V triangles
$T_2,T_3,T_4,T_5$ nonsupercritical Vd0.  Then the perimeter together with
$r_2$ is not covered.  The reflected $T_5$ placement is also impossible.
\end{lemma}

\begin{proof}
Let
\[
 I_L=\left[\frac{k}{W},R+\alpha\right],
 \qquad
 I_R=\left[\frac{k}{R},W+\delta\right].
\]
For the actual supercritical reaches $(A_0,B_0)$, put
\[
 A=\mathcal R_{I_R}(B_0),
 \qquad
 h_0=\mathcal R_{I_L}(A_0).
\]
The endpoint-distance inequality is
\[
 A_0^2+A_0B_0+B_0^2\le1.
\]
It implies $A_0,B_0<1$.  The signed center bounds give
$R+\alpha<1$, so $h_0>0$. Choose selected lower bounds
\[
 0\le a_i\le A_i,\qquad 0\le b_i\le B_i\qquad(1\le i\le5)
\]
so that boundary propagation gives
\[
 a_1\ge A,
 \qquad
 b_i\ge h_0\quad(1\le i\le5).
\]
The Vd cap gives
\begin{equation}
 A+h_0<\frac12.
 \label{eq:signed-adjacent-AH}
\end{equation}

The common perimeter budget gives
\[
 T:=\alpha+\delta<\frac1{24},
 \qquad
 T<\frac{\min\{R,W\}}6.
\]
We claim
\begin{equation}
 \delta<\frac {h_0}4.
 \label{eq:signed-adjacent-quarter}
\end{equation}
If $h_0=W-\alpha$, then
\[
 4\delta+\alpha\le4T<\frac{2W}{3}<W.
\]
Suppose $h_0=1-A_0$.  The other residual cannot be $1-B_0$, since otherwise
\eqref{eq:signed-adjacent-AH} would give $A_0+B_0>3/2$, whereas
the endpoint-distance inequality gives
$A_0+B_0\le2/\sqrt3<3/2$.  Hence
\[
 A=1-(W+\delta),
 \qquad
 B_0\ge y:=\frac{k}{R},
\]
and $W+\delta>1/2+h_0$.  Diameter transfer gives
\[
 h_0>\frac y2>\frac{\eta}{2R}.
\]
Thus
\[
 \delta>\vartheta(R):=R-\frac12+\frac{\eta}{2R}.
\]
Writing $\eta/R=W/(1+E)$, exact differentiation gives
\[
 4E(1+E)^2\vartheta'(R)
 =2E(1+E)(1+2E)-W(2R-1).
\]
If $R\le1/2$, the second term is nonnegative after subtraction.  If
$R\ge1/2$, then $W(2R-1)\le1/8$, while the first term is larger
than $1/8$.  Thus $\vartheta'(R)>0$.  Since $\vartheta(3/8)=1/24$, the bound
$\delta<1/24$ forces $R<3/8$.  Then
\[
 (2-3R)^2-E^2=(3-8R)(1-R)>0,
\]
so $\eta/R>1/3$ and $h_0>1/6$.  Hence $4\delta<1/6<h_0$, proving
\eqref{eq:signed-adjacent-quarter}.

The center interval on $r_2$ begins from the vertex side at $1-\delta$.  If
$T_1$ supports $r_2$, its exact endpoint satisfies
\[
 u_2<1-b_1-\frac{a_1}{t}\le1-h_0<1-\delta.
\]
Thus it does not bridge to the center.  For $T_2$, boundary handoff gives
\[
 a_2>\frac12+A,
 \qquad b_2\ge h_0.
\]
Put $p=1/2+A$.  Equation~\eqref{eq:signed-adjacent-AH} gives
$0\le h_0\le p$ and $p+h_0<1$.  Lemma~\ref{lem:quarter-radial-envelope} gives
\[
 c_2\le c_{\max}(p,h_0)
 \le1-\frac {h_0}4
 <1-\delta.
\]
Thus neither possible radial bridge reaches the center interval, and $r_2$ is
uncovered.
\end{proof}

\subsection{Nonadjacent Vd1/Vd2 placements}

\begin{lemma}[Nonadjacent Vd1/Vd2 placements]
\label{lem:signed-vd-nonadjacent-placement}
\label{thm:s2-geom-vd-nonadjacent}
Under the same reduced CE2 branch, let the unique Vd1/Vd2 V triangle be
$T_\tau$ with $\tau\in\{2,3,4\}$.  Then the perimeter together with
$r_\tau$ is not covered.
\end{lemma}

\begin{proof}
Put
\[
 T=\alpha+\delta,
 \qquad
 x=\frac{\eta+T}{W},
 \qquad
 y=\frac{\eta+T}{R}.
\]
Lemma~\ref{lem:signed-endpoints-dominate-slack} gives
\[
 T<\frac12\min\{x,y\}.
\]
Let $(A_0,B_0)$ be the actual reaches of the unique supercritical V triangle.
Then
\[
 A_0+B_0>1,\qquad A_0^2+A_0B_0+B_0^2\le1.
\]
Let
\[
 B=e_{[y,W+\delta]}(B_0),
 \qquad U=e_{[x,R+\alpha]}(A_0),
 \qquad A=1-B,
 \qquad h_0=1-U.
\]
Choose $a_\tau\le A_\tau$ and $b_\tau\le B_\tau$.  Boundary propagation
and the Vd cap permit the choices
\[
 a_\tau\ge A,
 \qquad b_\tau\ge h_0,
 \qquad A+h_0<\frac12.
\]

The component endpoint $B$ is either
$B_0$ or $W+\delta$.  In the latter case the center diameter gives
$B\le M_0(x)$ by Lemma~\ref{lem:signed-diameter-transfer}.  In the former
case, if $A_0<x$, then $U=A_0$ and
$A+h_0<1/2$ would force $A_0+B_0>3/2$, contrary to
$A_0+B_0\le2/\sqrt3$.  Hence $A_0\ge x$, and
$B=B_0\le M_0(A_0)\le M_0(x)$, using the decreasingness in that lemma.
Reflection gives $U\le M_0(y)$.
Consequently
\[
 A>\frac x2>T,
 \qquad
 h_0>\frac y2>T.
\]
Thus
\[
 d_\tau^C<\min\{A,h_0\}
 \qquad(\tau=2,3,4),
\]
because the three exits are $\delta$, $\alpha$, and
$\min\{\alpha/R,\delta/W\}\le T$.

The Vd corner normal form gives the radial margin
\[
 c_\tau<1-\min\{A,h_0\}.
\]
But the center interval begins from the vertex side at
\[
 1-d_\tau^C>1-\min\{A,h_0\}.
\]
The two radial traces are disjoint, and all other triangles are excluded by Vd0
locality and diameter locality.
\end{proof}

\subsection{Placement interfaces}

The local profile and radial-separation lemmas above are assembled exactly
once in Section~\ref{sec:strategy2-placement-assembly}.  This subsection
contains no independent CE2 one-Vd placement assembly.
